\documentclass{article}
\usepackage[utf8]{inputenc}
\usepackage{amsmath}
\usepackage{amsfonts}
\usepackage{amssymb}

\begin{document}
	
	\title{ Geometría Diferencial en la Manufactura e Industria}
	\author{Prof. David Duarte}
	\date{\today}
	
	\maketitle
	
	\section*{Problema}
	
	Consideren una superficie \( S \) en \( \mathbb{R}^3 \) definida por una función vectorial \( \vec{r}(u, v) \). Supongan que esta superficie se va a utilizar como molde en un proceso de manufactura. Utilizando los conceptos de geometría diferencial, como la primera y segunda forma fundamental, calculen y analicen las curvaturas principal, media y gaussiana en puntos específicos de \( S \). Luego, discutan cómo las propiedades de curvatura de \( S \) podrían influir en la selección de materiales y métodos de manufactura para minimizar el estrés y maximizar la durabilidad del producto final.
	
	\subsection*{Objetivos}
	
	\begin{enumerate}
		\item Aplicar la teoría de superficies para calcular curvaturas.
		\item Relacionar las curvaturas con consideraciones prácticas en el proceso de manufactura.
		\item Discutir la influencia de la geometría en la selección de materiales y métodos de producción.
		\item Modelar en el lenguaje de programación de su preferencia, un ejemplo tangible.
	\end{enumerate}
	
	\subsection*{Discusión}
	\begin{itemize}
		\item Discuta cómo la geometría diferencial proporciona el lenguaje matemático para formular y resolver problemas de la vida real, en el ámbito industrial.
		\item Reflexione sobre la importancia de comprender fenómenos asociados al comportamiento de los materiales y su aplicación en el campo industrial.
	\end{itemize}
	
		\section*{Entrega}
	Prepare un informe detallado que incluya todos los cálculos, diagramas y discusiones necesarias para cada uno de los puntos anteriores. El informe debe demostrar comprensión teórica y habilidad para aplicar conceptos matemáticos a problemas físicos.
	
\end{document}
